\subsection{The Odd Root Property} 
We can use radicals to find the solutions of equations, by rewriting power functions in terms of radical functions.
\begin{definition}[Odd Root Property]
For any expression $X$, and real number $c$, and odd integer $p$, the equations 
$$X^p=c\quad\mbox{and}\quad X=\sqrt[p]{c}$$
 are equivalent.
\end{definition}
\begin{Note}This is a case of a general move. If $f$ is a one-to-one function and $c$ is a number in the range of $f$, the equations $f(X)=c$ and $X=f\inv(c)$ are equivalent.\end{Note}
\begin{procedure}{Solving Equations with the Odd Root Property}
\begin{enumerate*}
\item Isolate the power expression.
\item Apply the odd root property to rewrite.
\item Continue solving.
\end{enumerate*}
\end{procedure}
\begin{Example}
Solve each equation below using the Odd Root Property.
\begin{lpicvsplit}{3}{7}
\foreach \x/\eqn in {0/x^3=125,1/(2x+1)^5=32,2/4(x-3)^7+8=4} {
	\regtext{\x*\value{fcols}+\x+\value{fcols}/2}{-1}{$\eqn$};
	\regtextleft{\x*\value{fcols}+\x}{-7}{Check:};
}
\end{lpicvsplit}
\end{Example}
We can solve some equations of this type as polynomial equations.
\begin{Example}
Solve $x^3=125$ as a polynomial equation by factoring.
\begin{lpic}{}{5}
\end{lpic}
\end{Example}